\documentclass[a4paper,14pt]{article}
\usepackage[utf8]{inputenc}
\usepackage[russian]{babel}
\usepackage[left=25mm,right=20mm,top=20mm,bottom=20mm]{geometry}
\usepackage{setspace}
\usepackage{indentfirst}
\usepackage{ulem} % Для корректного подчеркивания длинных строк

\begin{document}
\thispagestyle{empty} % Отключаем нумерацию страниц для титульного листа
\onehalfspacing

% Шапка организации
\begin{center}
    \small\bfseries
    ФЕДЕРАЛЬНОЕ ГОСУДАРСТВЕННОЕ АВТОНОМНОЕ ОБРАЗОВАТЕЛЬНОЕ \\
    УЧРЕЖДЕНИЕ ВЫСШЕГО ОБРАЗОВАНИЯ \\
    «САНКТ-ПЕТЕРБУРГСКИЙ ПОЛИТЕХНИЧЕСКИЙ УНИВЕРСИТЕТ \\
    ПЕТРА ВЕЛИКОГО» \\
    ФИЗИКО-МЕХАНИЧЕСКИЙ ИНСТИТУТ \\
    ВЫСШАЯ ШКОЛА ПРИКЛАДНОЙ МАТЕМАТИКИ И ВЫЧИСЛИТЕЛЬНОЙ ФИЗИКИ
\end{center}

\vspace{45mm}

% Название документа и тема
\begin{center}
    \large\bfseries
    Отчет о прохождении преддипломной практики \\[0.5em]
    на тему: «ОПТИМИЗАЦИЯ ЭВОЛЮЦИОННЫМ АЛГОРИТМОМ МЕТОДА РЕШЕНИЯ ЗАДАЧИ КОШИ НА БАЗЕ МЕТОДА РУНГЕ-КУТТЫ ДЛЯ ОСЦИЛЛИРУЮЩЕЙ ФУНКЦИИ»
\end{center}

\vspace{10mm}

% ФИО студента и группа
\begin{center}
    \uline{Добрецовой Елизаветы Викторовны, гр. 5040102/40101}
\end{center}

\vspace{25mm}

% Информационный блок руководителя и параметров
\noindent \textbf{Направление подготовки:} \uline{01.04.02 Прикладная математика и информатика}. \\
\noindent \textbf{Место прохождения практики:} \uline{СПбПУ Петра Великого}. \\
\noindent \textbf{Сроки практики:} \uline{с 17.04.2026 по 01.06.2026}. \\
\noindent \textbf{Руководитель практики от ФГАОУ ВО «СПбПУ»:} \uline{Козлов Константин Николаевич, доцент ВШПМиВФ, к.б.н.} \\
\noindent \textbf{Оценка:} \uline{\hspace{40mm}}

\vfill

% Подписи сторон внизу листа
\noindent
\begin{tabular*}{\textwidth}{@{\extracolsep{\fill}}l r}
    Руководитель практики & \\
    от ФГАОУ ВО «СПбПУ» & К.Н. Козлов \\[2em]
    Обучающийся & Е.В. Добрецова \\[2em]
    Дата: \uline{01.06.2026} & 
\end{tabular*}

\end{document}